\begin{table}[htbp]
\centering
\caption{Two-Way FE Decomposition---Table 2: Road Construction (Burgess et al., 2015)}
\label{tab:burgess_twfe_t2}
\begin{tabular}{lcc}
\toprule
 & FE (feTR) & FD (fdTR) \\
\midrule
\multicolumn{3}{l}{\textit{Panel A: Specification}} \\[3pt]
Regression type & Fixed Effects & First Differences \\
Dependent variable & \multicolumn{2}{c}{Change in paved road share} \\
Treatment variable & \multicolumn{2}{c}{President's coethnic (non-binary)} \\
Panel & \multicolumn{2}{c}{Districts $\times$ years (1964--2002)} \\[6pt]
\multicolumn{3}{l}{\textit{Panel B: Weight Decomposition}} \\[3pt]
$\hat{\beta}_{TWFE}$ & 1.911931 & --- \\
\# positive weights & 65 & --- \\
\# negative weights & 0 & --- \\
\% negative weights &   0.0\% & ---\% \\[6pt]
\multicolumn{3}{l}{\textit{Panel C: Classification (dCDH Web Appendix)}} \\[3pt]
Design & \multicolumn{2}{c}{Not binary and/or not staggered} \\
Dynamic effects & \multicolumn{2}{c}{No} \\
\bottomrule
\end{tabular}

\vspace{6pt}
\begin{minipage}{0.92\textwidth}
\footnotesize
\textit{Notes:} Same treatment variable and decomposition as Table 1, 
but using road construction data from maps (1964--2002) instead of 
road expenditure data (1963--2011). 
Negative weights indicate that the TWFE coefficient may not recover 
a convex combination of causal effects under heterogeneous treatment effects.
\end{minipage}
\end{table}
